\begin{trapbox}

	\begin{description}[style=nextline, leftmargin=2.8em, labelsep=0.5em, itemsep=0.35em]
		\item[\textbf{\textcolor{CTrap}{易错点一（SSA误用）}}]
			错误认为$AD=ED,\ CD=BD,\ \angle BAD=\angle CED$可证全等。
		\item[\textbf{\textcolor{CTrap}{正确理解（SAS判定）：}}]
			必须用$AD=ED,\ \angle ADC=\angle EDB,\ CD=BD$构成SAS。
		\item[\textbf{\textcolor{CTrap}{错因分析（全等混淆）：}}]
			混淆全等判定条件，SSA不成立。
	\end{description}

	\medskip
	\begin{description}[style=nextline, leftmargin=2.8em, labelsep=0.5em, itemsep=0.35em]
		\item[\textbf{\textcolor{CTrap}{易错点二（等长推平行）}}]
			由$AC=BE$直接断言$AC\parallel BE$，未通过角相等证明。
		\item[\textbf{\textcolor{CTrap}{正确理解（角条件平行）：}}]
			需由$\angle ACD=\angle EBD$得到平行。
		\item[\textbf{\textcolor{CTrap}{错因分析（直观误导）：}}]
			误以为等长即平行。
	\end{description}

	\medskip
	\begin{description}[style=nextline, leftmargin=2.8em, labelsep=0.5em, itemsep=0.35em]
		\item[\textbf{\textcolor{CTrap}{易错点三（忽略对称）}}]
			只知连接$BE$，不知连接$CE$同样可得结论。
		\item[\textbf{\textcolor{CTrap}{正确理解（对称辅助线）：}}]
			连接$CE$可得$AB=CE$且$AB\parallel CE$，两个方向都可。
		\item[\textbf{\textcolor{CTrap}{错因分析（思维定式）：}}]
			缺乏对辅助线多种选择的尝试。
	\end{description}

\end{trapbox}
